By the corrected construction in \(\texttt{def:binary-bow-menu}\), the three treatment-side branch effects are
\[
 \phi_0=\tfrac12 e_0,\qquad
 \phi_1=\tfrac12 e_1,\qquad
 \phi_2=\tfrac12\mathbf 1,
\]
and their associated states are \(\psi^{+},\psi^{+},\psi^{-}\), respectively.  Since \(e_0+e_1=\mathbf 1\),
\[
 \phi_0+\phi_1+\phi_2
 =\tfrac12(e_0+e_1+\mathbf 1)
 =\mathbf 1.
\]
Thus the unique treatment instrument obeys \(\texttt{ass:instrument-normalization}\).  The terminal effects obey \(e_0+e_1=\mathbf 1\), so the retained terminal basis instrument is normalized as well.  All displayed effects are nonnegative, and \(\psi^{+},\psi^{-},e_0,e_1\) are probability vectors by their definitions, so every listed pair is a well-typed separable branch.

The two treatment effects \(\tfrac12e_0\) and \(\tfrac12e_1\) are linearly independent, hence the treatment effect span is \(\mathbb R^2\).  Moreover
\[
 \det\!\begin{pmatrix}3/4&1/4\\[2pt]1/4&3/4\end{pmatrix}
 =\frac{9-1}{16}=\frac12\ne0,
\]
so \(\psi^{+}\) and \(\psi^{-}\) are linearly independent and the treatment state span is \(\mathbb R^2\).  At the terminal vertex the effects \(e_0,e_1\) form the standard basis, and for \(i\in\{0,1\}\),
\[
 (e_i\otimes e_i)(a_{\mathrm t},z_{\mathrm t})
 =\mathbf 1\{a_{\mathrm t}=i\}\mathbf 1\{z_{\mathrm t}=i\}
 =d_i(a_{\mathrm t},z_{\mathrm t}).
\]
In particular, the terminal effect-span condition in \(\texttt{thm:bow-rowspan-iff}\) holds for both outcome values.

By \(\texttt{def:bow-probe-matrix}\) and the corrected menu, the complete list of treatment rows is
\[
 w_0=\tfrac12e_0\otimes\psi^{+},\qquad
 w_1=\tfrac12e_1\otimes\psi^{+},\qquad
 w_2=\tfrac12\mathbf 1\otimes\psi^{-}.
\]
Suppose \(\alpha w_0+\beta w_1+\gamma w_2=0\).  Restricting this equality to the two coordinates with \(a_{\mathrm b}=0\) gives
\[
 \tfrac12\alpha\psi^{+}+\tfrac12\gamma\psi^{-}=0.
\]
The already established independence of \(\psi^{+},\psi^{-}\) implies \(\alpha=\gamma=0\).  Restriction to the coordinates with \(a_{\mathrm b}=1\) then gives \(\tfrac12\beta\psi^{+}=0\); since \(\psi^{+}\ne0\), \(\beta=0\).  Hence \(w_0,w_1,w_2\) are linearly independent and
\[
 \operatorname{rank}(A_{\mathrm b}^{\dagger})=3<4.
\]

The target-row identities are exact.  First,
\[
 w_0+w_1
 =\tfrac12(e_0+e_1)\otimes\psi^{+}
 =\tfrac12\mathbf 1\otimes\psi^{+}.
\]
Using \(\psi^{+}=(3/4,1/4)\) and \(\psi^{-}=(1/4,3/4)\),
\[
 \frac{3\psi^{+}-\psi^{-}}2=e_0,
 \qquad
 \frac{-\psi^{+}+3\psi^{-}}2=e_1.
\]
Therefore, by the definitions \(r_i=\mathbf 1\otimes e_i\),
\[
 3(w_0+w_1)-w_2
 =\mathbf 1\otimes\frac{3\psi^{+}-\psi^{-}}2
 =r_0,
\]
and
\[
 -(w_0+w_1)+3w_2
 =\mathbf 1\otimes\frac{-\psi^{+}+3\psi^{-}}2
 =r_1.
\]
Thus \(r_0,r_1\in\operatorname{row}(A_{\mathrm b}^{\dagger})\).  For every \(z^{\circ},y\in\{0,1\}\), both the treatment row-span condition and the terminal effect-span condition of the already-established \(\texttt{thm:bow-rowspan-iff}\) hold.  That theorem therefore yields
\[
 \operatorname{GID}_{\mathcal I_{\mathrm b}^{\dagger}}(Q_{z^{\circ},y}).
\]
Its global predicate ranges over the primitive-positive canonical edge-source class by \(\texttt{def:global-probe-identification}\), so this conclusion has exactly the claimed latent-bow scope.  Finally, the treatment split-port table has four coordinates whereas its probe matrix has rank three; hence the corresponding linear observation map has a nonzero nullspace and cannot tomographically recover the full treatment table.  This proves query-specific identification without full treatment-table tomography.